% Generated draft: requires mathematical and source review.
\documentclass[11pt]{article}
\usepackage[margin=25mm]{geometry}
\usepackage{fontspec}
\setmainfont{Latin Modern Roman}
\usepackage{amsmath,amssymb,mathtools,unicode-math}
\setmathfont{Latin Modern Math}
\usepackage{enumitem,needspace,xurl,hyperref}
\setlength{\parindent}{0pt}
\setlength{\parskip}{5pt}
\setlength{\emergencystretch}{3em}
\newcommand{\N}{\mathbb{N}}
\newcommand{\Z}{\mathbb{Z}}
\newcommand{\Q}{\mathbb{Q}}
\newcommand{\R}{\mathbb{R}}
\newcommand{\C}{\mathbb{C}}
\newcommand{\F}{\mathbb{F}}
\newcommand{\A}{\mathbb{A}}
\newcommand{\PP}{\mathbb{P}}
\newcommand{\E}{\mathbb{E}}
\newcommand{\eps}{\varepsilon}
\newcommand{\dd}{\,\mathrm{d}}
\begin{document}
\section*{UnsolvedMath ID: 30001135}
\textbf{Source problem code:} catalogue entry
\subsection{Definitions and mathematical statement}
Let \(K_4\) denote the complete graph on four vertices. A wheel graph is obtained from a cycle by adjoining a hub vertex adjacent to every cycle vertex. The supplied record does not specify the permitted cycle lengths or geometric degeneracies.

For terminology, in the usual real Euclidean framework convention, a planar framework consists of a finite graph \(G=(V,E)\) and positions \(p_i\in\R^2\), \(i\in V\). A self-stress is an assignment \(\omega:E\to\R\), with \(\omega_{ij}=\omega_{ji}\), satisfying the equilibrium equations
\[
\sum_{j:\{i,j\}\in E}\omega_{ij}(p_j-p_i)=0
\qquad(i\in V).
\]
This definition supplies conventional terminology only: the record does not specify whether these are precisely the conventions used in the referenced identities, nor whether ``planar'' requires a crossing-free drawing.

\textbf{Source question, with unresolved qualifications.} For the self-stress identities on planar \(K_4\) and wheel graphs described in the cited 2009 report, including the displayed relation
\[
\sum_{i<j}\omega_{ij}f_{ij}=1,
\]
seek a general mathematical construction explaining these identities and their connections to homology and cohomology. In particular, determine whether a homological or cohomological construction accounts for them.

The displayed relation is reproduced as a source-dependent formula, not asserted for every framework or every self-stress. The supplied record does not define \(f_{ij}\), its domain or codomain, the indexing set of the sum, the stress normalization, or the admissible geometric configurations. It also does not reproduce the other identities to be explained. Consequently, no precise universal quantifiers, chain complex, coefficient system, or proposed homological correspondence can be specified from this record. The question remains a research request rather than a fully specified proposition; these missing data require source consultation and human mathematical review.
\subsection{Short English statement}
The question asks whether the reported self-stress identities for planar \(K_4\) and wheel graphs arise from a broader construction, and how homology or cohomology explains them. The supplied record leaves the quantities \(f_{ij}\), normalization, geometric hypotheses, and additional identities unspecified, so it does not determine a precise theorem to prove or disprove.
\subsection{Sources}
\begin{itemize}
\item Local registry: \texttt{data/unsolvedmath/problems.json}, supplied record for permanent UnsolvedMath ID \(30001135\), entitled ``Homological Structure of Planar Self-Stress Identities,'' source problem code \texttt{catalogue entry}, in the Oberwolfach Reports Open Problems collection.
\item Supplied report citation: \emph{Discrete Differential Geometry} (2009), \emph{Oberwolfach Reports}, pp.~75--144, DOI \texttt{10.4171/owr/2009/02}. \url{https://doi.org/10.4171/owr/2009/02}. This is also the supplied literature-source link labelled ``original OWR report.'' The report's relevant passage is not reproduced in the supplied record.
\item Supplied literature source: Rote, \emph{Open Problems in Discrete Differential Geometry}. \url{https://page.mi.fu-berlin.de/rote/Kram/OWR-DDG09-problems.pdf}. The record supplies no named proposer; this reference does not by itself establish proposer attribution.
\item Supplied literature assessment dated \texttt{2026-08-21}: the status is labelled ``open,'' with the qualification that the extracted question broadly asks for connections to homology and cohomology and supplies no specific theorem whose truth or falsity can be checked. The background describes this as dated literature triage derived from the supplied Oberwolfach curation and requests independent verification of time-sensitive status and cited claims. This is supplied dated curation, not an independently verified assessment; neither linked source was visited for this draft.
\item The registry labels the statement \texttt{corrected\_verified} and describes its editing as meaning-preserving semantic, typographic, and LaTeX cleanup without unsupported source-dependent additions. These are supplied editorial labels, not independent verification of the mathematical formulation or its open status.
\end{itemize}
\end{document}
